%1/2in left/right margins, 1in top, 1/2 bottom
\documentclass[12pt]{article}
\hbadness=10000
\tolerance=10000
\textheight 9.7in
\textwidth 7.0in
\oddsidemargin -.5in
\topmargin -.5in
\headsep 0in
\headheight 0in
%\pagestyle{empty}

\newcommand{\by}{\mbox{\boldmath$y$}}
\newcommand{\bX}{\mbox{\boldmath$X$}}
\newcommand{\bbeta}{\mbox{\boldmath$\beta$}}

\begin{document}
\hrule
\medskip
\centerline{\textbf{STATISTICS 801 - Section 250}}\hfill \textbf{Fall 2016}
\centerline{\textbf{Why is this the test statistics for testing hypotheses for proportions?}}
\smallskip
\hrule
For a binomial random variable $E(X)=np$ and $V(X)=np(1-p)$.\\
Test statistic for proportions:\\
\begin{equation}
z=\frac{\hat{p}-p}{\sqrt{p(1-p)/n}}
\end{equation}
where $\hat{p}=\frac{X}{n}$ (sample proportion),\\
p is the population proportion, and n is the sample size.\\
The formula above for the test statistics can be derived from the $z=\frac{X-\mu}{\sigma}$ formula. Since $\mu=np$ and $\sigma=\sqrt{np(1-p)}$, we can write the following:\\
$z=\frac{X-np}{\sqrt{np(1-p)}}=\frac{X/n-np/n}{\sqrt{np(1-p)}/n}=\frac{\hat{p}-p}{\sqrt{np(1-p)/n^2}}=\frac{\hat{p}-p}{\sqrt{p(1-p)/n}}$.
\end{document}

